\begin{abstract}
First Jobbers must divide a limited monthly surplus between immediate liquidity, an emergency buffer, and a short-term financial goal. A fixed saving rule ignores the fact that future essential expenses and income volatility differ across users. We formulate this decision as a six-cycle finite-horizon Markov decision process with hard allocation feasibility and an evaluated shortfall-risk constraint. The proposed K PLUS experience recommends an allocation at the payday moment while leaving the final decision with the user. The study compares a transparent fixed rule, known-model Dynamic Programming (DP), Bayesian adaptive DP, pooled Reinforcement Learning (RL), and a small recurrent Meta-RL policy. The recurrent policy receives no task label and updates its context from observed salary, expense, and reward history. All methods use the same simulator, action grid, reward, and matched evaluation paths. The primary research question is whether history-based Meta-RL improves held-out cash-flow adaptation enough to justify its complexity over Bayesian DP. The work is an auditable Track 2 experiment rather than a production banking integration or a claim about real customer outcomes.
\end{abstract}

\begin{IEEEkeywords}
First Jobber, sequential decision-making, Dynamic Programming, Bayesian adaptation, Meta-Reinforcement Learning, liquidity risk, financial goal allocation
\end{IEEEkeywords}
